\chapter{Dental Braces and Retainers}

\section*{Definition and Overview}
Dental braces are orthodontic appliances used to straighten teeth, correct bites, and align the jaws. They work by applying gentle, continuous pressure that moves teeth through the bone over time. Braces may be fixed (bonded to the teeth) or removable, including clear aligners. After active treatment, retainers are worn to hold the teeth in their new positions, because teeth tend to drift back toward their original alignment. This chapter explains the types of braces, how treatment works, and how to care for teeth and appliances during and after orthodontic treatment.

\section*{Epidemiology}
Orthodontic treatment is common, with millions of people worldwide wearing braces at any time. In Australia, a significant proportion of children and teenagers, and increasing numbers of adults, receive orthodontic care. Around one in three adults feels their teeth could be straighter, and orthodontic treatment is often undertaken for functional reasons such as bite problems, as well as cosmetic ones. Treatment typically lasts 12--24 months for fixed braces, with clear aligner therapy offering a popular alternative, especially for adults.

\section*{Aetiology and Pathophysiology}
\begin{itemize}
    \item \textbf{Malocclusion:} misalignment of the teeth or jaws, including crowding, gaps, overbites, underbites, and crossbites
    \item \textbf{Genetics:} jaw size, tooth size, and bite relationships are inherited
    \item \textbf{Habits:} thumb sucking, tongue thrusting, and prolonged dummy use can affect dental alignment
    \item \textbf{Tooth loss:} early loss of baby teeth can allow permanent teeth to drift
    \item \textbf{Mechanism of braces:} brackets and wires apply forces that stimulate bone remodelling, moving teeth through the jawbone
    \item \textbf{Retention:} after teeth are moved, new bone must stabilise, and retainers prevent relapse to the original position
\end{itemize}

\section*{Clinical Features}
\begin{itemize}
    \item Crooked, crowded, or gapped teeth
    \item Overbite, underbite, or open bite affecting the bite and appearance
    \item Difficulty cleaning between misaligned teeth, increasing decay risk
    \item Speech difficulties or jaw pain in some bite problems
    \item With braces: initial discomfort and mouth ulcers for the first days after adjustments
    \item Loose brackets, broken wires, or sore spots during treatment
    \item After braces: teeth may feel slightly loose as they move; retainer wear prevents them drifting back
\end{itemize}

\section*{Differential Diagnosis}
\begin{itemize}
    \item Malocclusion requiring orthodontic treatment versus mild misalignment managed by observation
    \item Jaw joint (TMJ) disorders, which can be associated with bite problems
    \item Tooth decay or gum disease, which may coexist and need treatment before or during orthodontics
    \item Missing or impacted teeth, which may require surgical or restorative planning
    \item Sleep-disordered breathing or mouth breathing, which can affect dental and jaw development
\end{itemize}

\section*{When to Seek Medical Care}
Children should have an orthodontic assessment around age 7--8 if bite or alignment problems are noticed, and most treatment begins in the early teenage years. Adults can be treated at any age if gums and bone are healthy. See the orthodontist promptly for broken brackets, loose wires, or significant pain, and seek urgent care for trauma to braced teeth.

\textbf{Call 000 (or local emergency services) for:}
\begin{itemize}
    \item Trauma to the face or teeth, including a knocked-out or fractured tooth with braces
    \item A wire that is causing severe injury to the gums, cheeks, or tongue
    \item Swelling or signs of infection around a tooth
\end{itemize}

\section*{Diagnosis and Investigations}
\begin{itemize}
    \item \textbf{Orthodontic examination:} assessment of teeth, bite, and facial profile
    \item \textbf{Dental X-rays:} panoramic and bitewing X-rays assess tooth position, roots, and developing teeth
    \item \textbf{Photographs and impressions:} records used to plan treatment; digital scans may replace impressions
    \item \textbf{Cephalometric analysis:} X-ray of the skull used to assess jaw relationships
    \item \textbf{Treatment planning:} the orthodontist plans the appliance type and treatment sequence
\end{itemize}

\section*{Management}
\begin{itemize}
    \item \textbf{Fixed braces:} metal or ceramic brackets bonded to the teeth, with archwires adjusted regularly
    \item \textbf{Clear aligners:} a series of removable transparent trays, such as Invisalign, that progressively move teeth
    \item \textbf{Functional appliances:} removable devices that guide jaw growth in growing children
    \item \textbf{Headgear or elastics:} additional forces to correct severe bite problems
    \item \textbf{Retention:} fixed (bonded) or removable retainers worn after braces, usually full-time initially then at night
    \item \textbf{Regular adjustments:} appointments every 4--8 weeks during active treatment
    \item \textbf{Managing discomfort:} soft foods, pain relief, and orthodontic wax for mouth ulcers
    \item \textbf{Repairs:} prompt attention to broken brackets and wires
\end{itemize}

\section*{Complications}
\begin{itemize}
    \item Tooth decay and white spot lesions from poor cleaning around braces
    \item Gum inflammation and swelling from plaque retention
    \item Root resorption, a shortening of tooth roots in a minority of patients
    \item Relapse of tooth position if retainers are not worn as directed
    \item Prolonged treatment if appointments are missed or appliances are damaged
    \item Discomfort, mouth ulcers, and soft tissue irritation
\end{itemize}

\section*{Prognosis and Outcomes}
Orthodontic treatment is highly effective when appliances are worn and oral hygiene is maintained. Most people achieve well-aligned teeth and a corrected bite, with improvements in appearance, function, and confidence. Treatment results are permanent if retainers are worn as prescribed, since teeth have a strong tendency to relapse. The long-term outcome is excellent, and most people retain stable, healthy results with routine dental care.

\section*{Prevention and Self-Care}
\begin{itemize}
    \item Brush after every meal and floss or use interdental brushes daily during treatment
    \item Use fluoride toothpaste and consider additional fluoride treatments
    \item Avoid hard, sticky, and sugary foods that damage braces or increase decay
    \item Wear a mouthguard for sport
    \item Attend all orthodontic appointments
    \item Wear retainers exactly as instructed to prevent relapse
    \item Continue regular dental check-ups throughout treatment
\end{itemize}

\section*{Sources and Further Reading}
The following reputable sources provide further information for healthcare students and patients seeking reliable, current advice about dental braces and retainers.

\begin{itemize}
    \item Australian Society of Orthodontists. \textit{Braces and orthodontic treatment information.} https://www.aso.org.au/
    \item Australian Dental Association. \textit{Orthodontics: what to expect.} https://www.ada.org.au/
    \item NHS. \textit{Orthodontic treatment and braces.} https://www.nhs.uk/conditions/orthodontics/
    \item Mayo Clinic. \textit{Dental braces: types and care.} https://www.mayoclinic.org/
    \item MedlinePlus. \textit{Braces and orthodontics.} https://medlineplus.gov/
    \item Healthdirect Australia. \textit{Dental braces: information for children and adults.} https://www.healthdirect.gov.au/
\end{itemize}
